\section{Introduction}
\label{sec:intro}

Emergency communication is usually treated as a benign input to evacuation:
give people better information and their decisions should improve. That
assumption is plausible at the individual level, but it is not obviously true
at the crowd level. A warning that reduces uncertainty for everyone at once can
also synchronize route choice, push agents toward the same exit, amplify a
bottleneck, or move people through a contaminated region before the physical
hazard has dissipated. The question is therefore not simply whether information
helps, but when, where, and for whom it helps.

This paper studies evacuation as a coupled physical-information system. The
physical side includes layout geometry, social-force movement, bottleneck
dynamics, hazard spread, and physiological impairment. The information side
includes partial exit knowledge, hazard beliefs, signage and PA broadcasts,
responder credibility, local gossip, message decay, and distortion. Chiyoda
connects these layers by making agents plan through their believed world rather
than the ground truth world. They may avoid a hazard they have heard about,
miss an exit they do not know, herd when uncertainty is high, or follow a
stale route because the belief has not yet decayed.

\paragraph{Research question.} We ask whether entropy-guided communication
can improve evacuation outcomes without creating harmful convergence. More
concretely: which intervention policies reduce belief entropy and improve
belief accuracy per unit of induced exposure, queueing, and exit imbalance?
The current study scaffold compares static signage, global broadcasts,
responder relay, entropy-targeted broadcasts, density-aware targeting,
exposure-aware targeting, and bottleneck-avoidance targeting.

\paragraph{Contributions.} This paper makes four contributions.

\begin{enumerate}
\item \textbf{A controllable information-intervention layer.} Chiyoda turns
emergency communication into a first-class policy object rather than a fixed
scenario assumption. Policies can be static, authority-driven, or adaptive to
entropy, density, exposure, and bottleneck pressure.

\item \textbf{A coupled ITED simulation loop.} Agents carry probabilistic
beliefs over exits and hazards, update them through observation, beacon
broadcast, responder relay, and gossip, then route through belief-weighted
costs under social-force dynamics and hazard impairment.

\item \textbf{Information-safety metrics.} We introduce study outputs that
jointly measure entropy reduction, belief accuracy gain, intervention reach,
hazard exposure pressure, queue pressure, exit imbalance, and a harmful
convergence index.

\item \textbf{A reproducible paper pipeline.} YAML study definitions generate
multi-seed bundles with Parquet/CSV tables and paper figures. The LaTeX paper
ingests generated study statistics through \texttt{stats.tex}.
\end{enumerate}

\paragraph{Roadmap.} \S\ref{sec:background} positions the work against
pedestrian dynamics, information-aware evacuation, and ABM validation work.
\S\ref{sec:design} describes Chiyoda's model and intervention policies.
\S\ref{sec:implementation} explains the software pipeline. \S\ref{sec:evaluation}
lays out the planned empirical study. \S\ref{sec:related} sharpens the
distinction from prior work. \S\ref{sec:limitations} states the current
threats to validity, and \S\ref{sec:conclusion} concludes.
